\documentclass[11pt,a4paper]{article}
\usepackage[margin=2.5cm]{geometry}
\usepackage{amsmath,amssymb,amsthm,bm}
\usepackage[round]{natbib}
\usepackage[colorlinks,linkcolor=blue,citecolor=blue]{hyperref}

\newcommand{\bs}{\mathbf{s}}
\newcommand{\bh}{\mathbf{h}}
\newcommand{\bc}{\mathbf{c}}
\newcommand{\bT}{\mathbf{T}}
\newcommand{\bn}{\hat{\mathbf{n}}}
\newcommand{\br}{\hat{\mathbf{r}}}
\newcommand{\bnu}{\hat{\bm{\nu}}}
\newcommand{\bx}{\mathbf{x}}
\newcommand{\bM}{\mathbf{M}}
\newcommand{\vm}{\mathbf{m}}
\newcommand{\bI}{\mathbf{I}}
\newcommand{\sd}{\mathbf{s}^\circ}
\newcommand{\pd}{\phi^\circ}
\newcommand{\jump}[1]{\left[#1\right]^+_-}
\renewcommand{\Re}{\operatorname{Re}}
\newcommand{\sym}{\operatorname{sym}}
\newcommand{\dd}{\,\mathrm{d}}
\newcommand{\Lag}{\mathcal{L}}

\newtheorem{lemma}{Lemma}
\theoremstyle{remark}
\newtheorem{remark}{Remark}

\title{Adjoint kernels for a self-gravitating, hydrostatic earth model\\
from a weak formulation}
\author{Draft for checking}
\date{\today}

\begin{document}
\maketitle

\begin{abstract}
We derive sensitivity kernels for a non-rotating, self-gravitating earth
model in hydrostatic equilibrium, and restrict the perturbations to those that
keep it so: in three dimensions, perturbations of the elastic tensor
(including its anelastic dispersion) at fixed density, and of the source;
for spherically symmetric perturbations, also the density and the radii of
discontinuities. Laterally varying density and topography would introduce a
deviatoric initial stress that the hydrostatic equations do not determine,
and we show that equivalent forms of those equations then give different
first-order sensitivities. The derivation starts from the weak form of the
equations of motion, in which the boundary conditions are natural, and
treats the equilibrium gravitational potential $\Phi$ as an independent
variable constrained by Poisson's equation, so that the effect of a density
perturbation on the equilibrium gravity is carried by a single, static
adjoint potential $\lambda$ whose source is formed, once, from the forward
and adjoint fields integrated over frequency. Every kernel is then local, and
for spherically symmetric perturbations $\lambda$ reduces to a radial
quadrature. The method is formulated in the frequency domain at complex
frequencies, as in \citet{Adourian2024}; the adjoint problem is derived in
continuous form and then discretised with the rules that form the synthetic
seismograms, which for these linear operations coincides with the discrete
adjoint.
\end{abstract}

\section{Setting and notation}

The earth model occupies a ball $V$ of radius $a$, divided into regions by
internal discontinuities $\Sigma = \Sigma_{SS}\cup\Sigma_{FS}$ (solid--solid
and fluid--solid), with free surface $\partial V$. The unit normal $\bn$
points from the $-$ side of a boundary to the $+$ side, and
$\jump{f} = f_+ - f_-$; on $\partial V$ the $+$ side is outside. With this
convention, for a tensor field $\mathbf{v}$ smooth in each region,
\begin{equation}
  \int_V \nabla\cdot\mathbf{v}\dd V = -\int_{\Sigma\cup\partial V}
  \jump{\bn\cdot\mathbf{v}}\dd\Sigma, \label{eq:gauss}
\end{equation}
with $\mathbf{v}=0$ outside $V$. For vector fields we use the bilinear
pairing $\langle \mathbf{a},\mathbf{b}\rangle = \int_V\mathbf{a}\cdot
\mathbf{b}\dd V$, without complex conjugation; all forms below are bilinear,
and the frequency-domain operators are complex symmetric rather than
Hermitian.

The equilibrium model has density $\rho$, elastic tensor $\bc$ (the
incremental tensor of the prestressed state, with the major and minor
symmetries) and gravitational potential $\Phi$, with
\begin{equation}
  \nabla^2\Phi = 4\pi G\rho \quad\text{in } \mathbb{R}^3,
  \qquad \Phi\to0 \text{ as } |\bx|\to\infty. \label{eq:poisson-eq}
\end{equation}
It is in hydrostatic equilibrium: the initial stress is a pressure $p_0$
with $\nabla p_0 = -\rho\nabla\Phi$, so that $\nabla\rho\parallel\nabla\Phi$
in each region and the boundaries are level surfaces of $\Phi$. For a
spherically symmetric model, $\nabla\Phi = g\br$ with $g = 4\pi G
r^{-2}\int_0^r\rho r'^2\dd r' \ge 0$.

\paragraph{Scope.} The background is spherically symmetric. A perturbation of
$\bc$, or of its quality factors, at fixed $\rho$ leaves $\rho$, $\Phi$ and
$p_0$ unchanged, so the perturbed model is still in hydrostatic equilibrium,
whatever the perturbation's form; so does a spherically symmetric
perturbation of $\rho$ or of the radius of a discontinuity. These are the
perturbations considered: three-dimensional ones of the elastic tensor at
fixed density, and spherically symmetric ones of every parameter. A
laterally varying density or topography is excluded, for the reason given in
Section~\ref{sec:lateral}. The derivation below is written for general
perturbations of $\rho$ and $\Phi$ where that costs nothing, and specialised
at the end.

\section{The forward problem}

\subsection{The hydrostatic equations}

In the hydrostatic form \citep[][eq.~3.259]{DahlenTromp1998}, the initial
pressure has been eliminated with $\nabla p_0=-\rho\nabla\Phi$, and the
displacement $\bs$ and the perturbation $\phi$ of the gravitational
potential satisfy, without rotation,
\begin{align}
  \rho\ddot\bs - \nabla\cdot\bT + \nabla(\rho\bs\cdot\nabla\Phi)
    - \nabla\cdot(\rho\bs)\nabla\Phi + \rho\nabla\phi &= \mathbf{f},
  \qquad \bT = \bc:\nabla\bs, \label{eq:momentum}\\
  (4\pi G)^{-1}\nabla^2\phi + \nabla\cdot(\rho\bs) &= 0,
  \label{eq:poisson}
\end{align}
the first in each region of $V$, the second in $\mathbb{R}^3$. The boundary
conditions are $\jump{\bn\cdot\bT}=\mathbf{0}$ on $\Sigma\cup\partial V$
(with $\bT=0$ outside), $\jump{\bs}=\mathbf{0}$ on $\Sigma_{SS}$,
$\jump{\bn\cdot\bs}=0$ and $\bn\cdot\bT\parallel\bn$ on $\Sigma_{FS}$, and
$\jump{\phi}=0$, $\jump{\bn\cdot\bm{\xi}}=0$ with $\bm{\xi}=(4\pi
G)^{-1}\nabla\phi+\rho\bs$ on every boundary, with $\phi\to0$ at infinity
\citep{LiuTromp2008}. These are the equations that YSpec solves, written in
terms of the radial variables of Section~\ref{sec:yspec}.

Integrating \eqref{eq:momentum} through a thin layer across a boundary, the
distributional parts of $\nabla(\rho\bs\cdot\nabla\Phi)$ and
$-\nabla\cdot(\rho\bs)\nabla\Phi$ are $\bn\jump{\rho}g\,\bn\cdot\bs$ and
$-\jump{\rho}(\bn\cdot\bs)g\,\bn$, which cancel because $\nabla\Phi=g\bn$
there and $\bn\cdot\bs$ is continuous. The equations therefore hold in each
region with only the stated conditions on the boundaries.

\subsection{The weak form}

Let $\mathcal{U}$ be the displacements that honour the kinematic conditions
($\bs$ continuous on $\Sigma_{SS}$, $\bn\cdot\bs$ continuous on
$\Sigma_{FS}$) and $\mathcal{P}$ the potentials in $H^1(\mathbb{R}^3)$ that
decay at infinity. Take $\bs'\in\mathcal{U}$, multiply \eqref{eq:momentum}
by it, integrate over $V$, and use \eqref{eq:gauss}:
\begin{itemize}
\item $-\langle\bs',\nabla\cdot\bT\rangle = \int_V\nabla\bs':\bT\dd V +
  \int\jump{\bs'\cdot\bn\cdot\bT}\dd\Sigma$, and the surface term vanishes:
  on $\Sigma_{SS}$ both $\bs'$ and $\bn\cdot\bT$ are continuous, on
  $\Sigma_{FS}$ the traction is $\tau\bn$ with $\tau$ and $\bn\cdot\bs'$
  continuous, and on $\partial V$ the traction is zero.
\item $\langle\bs',\nabla(\rho\bs\cdot\nabla\Phi)\rangle
  -\langle\bs',\nabla\cdot(\rho\bs)\nabla\Phi\rangle
  = \int_V\rho\,[\bs\cdot\nabla(\bs'\cdot\nabla\Phi) - (\nabla\cdot\bs')
  (\bs\cdot\nabla\Phi)]\dd V$ plus the surface terms
  $\int\jump{-(\bn\cdot\bs')\rho\,\bs\cdot\nabla\Phi + (\bs'\cdot\nabla\Phi)
  \rho\,\bn\cdot\bs}\dd\Sigma$, which vanish because $\nabla\Phi=g\bn$ on
  the boundaries.
\end{itemize}
Similarly, multiplying \eqref{eq:poisson} by $\phi'\in\mathcal{P}$ and
integrating over $\mathbb{R}^3$, the surface terms vanish by the continuity
of $\bn\cdot\bm\xi$. With
\begin{align}
  A_g(\bs',\bs) &= \int_V\rho\left[\bs\cdot\nabla\nabla\Phi\cdot\bs'
    + (\bs\cdot\nabla\bs')\cdot\nabla\Phi - (\nabla\cdot\bs')(\bs\cdot
    \nabla\Phi)\right]\dd V, \label{eq:Ag}\\
  B\big((\bs',\phi'),(\bs,\phi)\big) &= \int_V\nabla\bs':\bc:\nabla\bs\dd V
    + A_g(\bs',\bs) \nonumber\\
    &\quad + \int_V\rho\,(\bs'\cdot\nabla\phi + \bs\cdot\nabla\phi')
    \dd V + \frac{1}{4\pi G}\int_{\mathbb{R}^3}\nabla\phi'\cdot\nabla\phi
    \dd V, \label{eq:B}
\end{align}
where $(\bs\cdot\nabla\bs')\cdot\nabla\Phi = s_i\,\partial_i s'_j\,
\partial_j\Phi$, the forward problem is to find $(\bs,\phi)$ in
$\mathcal{U}\times\mathcal{P}$ such that
\begin{equation}
  \langle\bs',\rho\ddot\bs\rangle + B\big((\bs',\phi'),(\bs,\phi)\big)
  = \langle\bs',\mathbf{f}\rangle
  \quad\text{for all } (\bs',\phi')\in\mathcal{U}\times\mathcal{P}.
  \label{eq:weak}
\end{equation}
The equation with $\phi'=0$ is the weak form of \eqref{eq:momentum}; that
with $\bs'=0$ is minus that of \eqref{eq:poisson}. All the boundary
conditions on tractions and fluxes are natural. The gravitational form $A_g$
is eq.~B26 of \citet{AlAttarTromp2014}, derived there for quasi-static
loading, where the fluid is described by $\phi$ alone and $A_g$ is accordingly
an integral over the solid with surface terms on the fluid--solid
boundaries; with a fluid displacement, as here, those become part of the
volume integral.

\begin{lemma}\label{lem:symmetry}
For a hydrostatic background, $A_g$, and so $B$, is symmetric.
\end{lemma}
\begin{proof}
With $\mathbf{N} = \rho(\bs\otimes\bs'-\bs'\otimes\bs)$,
\[
  A_g(\bs',\bs)-A_g(\bs,\bs') = \int_V\left[\partial_i N_{ij}\,
  \partial_j\Phi - (\partial_i\rho)(s_is'_j-s'_is_j)\partial_j\Phi\right]
  \dd V.
\]
The second term vanishes because $\nabla\rho\parallel\nabla\Phi$. In the
first, integrating by parts leaves $-\int N_{ij}\partial_i\partial_j\Phi$,
which vanishes because $\mathbf{N}$ is antisymmetric, and the surface term
$\int\jump{n_iN_{ij}\partial_j\Phi}\dd\Sigma$, which is
$\int\jump{g\,\bn\cdot\mathbf{N}\cdot\bn}\dd\Sigma=0$.
\end{proof}
The symmetry gives reciprocity, and with it the identification of the
adjoint problem with a forward one (Section~\ref{sec:adjoint}). It uses
the hydrostatic conditions; we return to this in
Section~\ref{sec:lateral}.

\subsection{Complex frequencies}

As in YSpec and \citet{Adourian2024}, the equations are transformed with
time dependence $e^{i\omega t}$ and solved at the complex frequencies
$\omega_\varepsilon = \omega - i\varepsilon$, $\omega$ real, with
$\varepsilon$ chosen so that the response decays by a factor $e^{-5}$ over
the length of the record; the time series is recovered by an inverse
transform followed by multiplication by $e^{\varepsilon t}$. For each real
$\omega$ the fields $(\bs_\omega,\phi_\omega)$, which we label by $\omega$
although they are evaluated at $\omega_\varepsilon$, satisfy
\begin{equation}
  S_\omega\big((\bs',\phi'),(\bs_\omega,\phi_\omega)\big) :=
  -\omega_\varepsilon^2\langle\bs',\rho\bs_\omega\rangle
  + B_\omega\big((\bs',\phi'),(\bs_\omega,\phi_\omega)\big) =
  \langle\bs',\mathbf{f}_\omega\rangle, \label{eq:S}
\end{equation}
where $B_\omega$ is $B$ with $\bc$ replaced by the complex tensor
$\bc_\omega=\bc(\omega_\varepsilon)$, which has the same symmetries. Each
$S_\omega$ is a complex-symmetric bilinear form. Because the time series are
real, the fields at $-\omega$ are the complex conjugates of those at
$\omega$, and for such Hermitian functions of frequency we write
\begin{equation}
  \langle\!\langle F\rangle\!\rangle = \frac{1}{2\pi}\int_{-\infty}^\infty
  F(\omega)\dd\omega = \frac{1}{\pi}\int_0^\infty\Re F(\omega)\dd\omega.
  \label{eq:fsum}
\end{equation}

\section{Observables and adjoint sources}\label{sec:observables}

\subsection{The continuous adjoint source}

Let $X(\omega) = \bnu\cdot\bs_\omega(\bx_r)$ be the spectrum of the
component $\bnu$ at the receiver $\bx_r$. The processed synthetic is
\begin{equation}
  u(t) = \frac{e^{\varepsilon t}}{2\pi}\int_{-\infty}^\infty w(\omega)
  R(\omega)X(\omega)\,e^{i\omega t}\dd\omega, \label{eq:synthetic}
\end{equation}
with $w$ a real taper and $R$ the instrument response at the complex
frequency, which includes any time derivatives of the displacement. For a
real functional $J[u]$ with Fr\'echet derivative $j(t) = \delta J/\delta
u(t)$,
\begin{equation}
  \delta J = \int j(t)\,\delta u(t)\dd t
  = \langle\!\langle A\,\delta X\rangle\!\rangle, \qquad
  A(\omega) = w(\omega)R(\omega)\int j(t)\,e^{\varepsilon t}e^{i\omega t}
  \dd t, \label{eq:A}
\end{equation}
the processed adjoint source of \citet[][eqs~37--38]{Adourian2024}: the
factor $e^{\varepsilon t}$ reflects the undamping of the synthetic.

\subsection{Cross-correlation traveltimes}

For the delay $\tau$ of a phase in a window $W(t)$, measured by
cross-correlation with a reference that is the synthetic itself
\citep{Dahlen2000}, a delay $\delta\tau$ changes the synthetic by
$\delta u = -\dot u\,\delta\tau$, and so
\begin{equation}
  \delta\tau = -\frac{\int W\dot u\,\delta u\dd t}{\int W\dot u^2\dd t},
  \qquad j(t) = -\frac{W(t)\dot u(t)}{\int W\dot u^2\dd t},
  \label{eq:traveltime}
\end{equation}
with $\dot u$ the velocity synthetic: the adjoint source of
\citet[][eqs~82--83]{LiuTromp2008}, here at complex frequency.

\section{The Lagrangian}\label{sec:adjoint}

The model parameters are $\rho$, $\bc$ (through its elastic values and
quality factors) and the source. The equilibrium potential depends on
$\rho$ through \eqref{eq:poisson-eq}; rather than eliminate it, we treat
$\Phi$ as an independent variable and impose \eqref{eq:poisson-eq} in weak
form with a multiplier $\lambda$:
\begin{equation}
  P(\Phi,\lambda;\rho) := \frac{1}{4\pi G}\int_{\mathbb{R}^3}\nabla\lambda
  \cdot\nabla\Phi\dd V + \int_V\rho\lambda\dd V = 0 \quad\text{for all }
  \lambda. \label{eq:P}
\end{equation}
The Lagrangian is
\begin{equation}
  \Lag = J - \left\langle\!\!\left\langle
  S_\omega\big((\sd_\omega,\pd_\omega),(\bs_\omega,\phi_\omega)\big)
  - \langle\sd_\omega,\mathbf{f}_\omega\rangle\right\rangle\!\!\right\rangle
  + P(\Phi,\lambda;\rho), \label{eq:L}
\end{equation}
with $(\sd_\omega,\pd_\omega)$ the multipliers for the forward problems,
Hermitian in $\omega$ like the fields. Where the forward and equilibrium
problems hold, $\Lag=J$ whatever the multipliers, so $\delta J =
\partial_m\Lag\,\delta m$ once the multipliers make $\Lag$ stationary with
respect to $\bs_\omega$, $\phi_\omega$ and $\Phi$.

\paragraph{The adjoint fields.}
Stationarity with respect to $(\bs_\omega,\phi_\omega)$ at each $\omega>0$,
using \eqref{eq:A}, requires $\Re[A\,\bnu\cdot\delta\bs(\bx_r)] =
\Re[S_\omega((\sd_\omega,\pd_\omega),(\delta\bs,\delta\phi))]$ for all
complex $(\delta\bs,\delta\phi)$, and so, replacing $\delta\bs$ by
$i\delta\bs$ for the imaginary part, and by the symmetry of $S_\omega$
(Lemma~\ref{lem:symmetry}),
\begin{equation}
  S_\omega\big((\delta\bs,\delta\phi),(\sd_\omega,\pd_\omega)\big) =
  \langle\delta\bs, A(\omega)\bnu\,\delta(\bx-\bx_r)\rangle
  \quad\text{for all } (\delta\bs,\delta\phi). \label{eq:adjoint}
\end{equation}
The adjoint field at $\omega$ is the forward response, at the same complex
frequency, to the point force $A(\omega)\bnu$ at the receiver. No time
reversal or change of frequency is needed, and the solver is the forward
one.

\paragraph{The adjoint potential.}
Stationarity with respect to $\Phi$ requires
\begin{equation}
  \frac{1}{4\pi G}\int_{\mathbb{R}^3}\nabla\lambda\cdot\nabla\psi\dd V =
  \int_V\left(\bM:\nabla\nabla\psi + \vm\cdot\nabla\psi\right)\dd V
  \quad\text{for all } \psi, \label{eq:lambda}
\end{equation}
where, from \eqref{eq:Ag}, the forward and adjoint fields enter through the
static fields
\begin{equation}
  \bM = \langle\!\langle\rho\,\sym(\bs_\omega\otimes\sd_\omega)
  \rangle\!\rangle, \qquad
  \vm = \langle\!\langle\rho\left((\bs_\omega\cdot\nabla)\sd_\omega -
  \bs_\omega\,\nabla\cdot\sd_\omega\right)\rangle\!\rangle.
  \label{eq:sources}
\end{equation}
This is a single, static Poisson problem, solved once after the frequency
integral; for the spherically symmetric density perturbations considered
here only its spherical average is needed, which is a radial quadrature
(Section~\ref{sec:lambda}).

\section{The kernels}

With the multipliers so chosen,
\begin{equation}
  \delta J = -\left\langle\!\!\left\langle\partial_mS_\omega
  \big((\sd_\omega,\pd_\omega),(\bs_\omega,\phi_\omega)\big)\delta m
  \right\rangle\!\!\right\rangle + \langle\!\langle\langle\sd_\omega,
  \delta\mathbf{f}_\omega\rangle\rangle\!\rangle + \int_V\lambda\,\delta\rho
  \dd V. \label{eq:gradient}
\end{equation}

\paragraph{Elastic tensor.}
$\delta J = \int_V \mathbf{K}_c::\delta\bc\dd V$ with
\begin{equation}
  \mathbf{K}_c = -\left\langle\!\!\left\langle\frac{\partial\bc_\omega}
  {\partial\bc}^{\!\mathsf T}\!\!:\big(\nabla\sd_\omega\otimes\nabla
  \bs_\omega\big)\right\rangle\!\!\right\rangle,
\end{equation}
the chain rule through the dispersion. For transverse isotropy with a radial
axis, in terms of the strains $\bm\epsilon=\sym\nabla\bs_\omega$ and
$\bm\epsilon^\circ=\sym\nabla\sd_\omega$, with $h$ the horizontal
components, the elastic energy density is
$C\epsilon^\circ_{rr}\epsilon_{rr} + A\,\mathrm{tr}_h\bm\epsilon^\circ\,
\mathrm{tr}_h\bm\epsilon + F(\epsilon^\circ_{rr}\mathrm{tr}_h\bm\epsilon +
\mathrm{tr}_h\bm\epsilon^\circ\epsilon_{rr}) + 4L\,\bm\epsilon^\circ_{rh}
\cdot\bm\epsilon_{rh} + N(2\bm\epsilon^\circ_{hh}:\bm\epsilon_{hh} -
2\,\mathrm{tr}_h\bm\epsilon^\circ\,\mathrm{tr}_h\bm\epsilon)$, every term
invariant under rotation about $\br$, so the two fields may be expressed in
different horizontal frames. With $D(\omega)$ the dispersion of YSpec, so
that a modulus $X$ with inverse quality factor $q_X$ is $X(1+q_XD)$ at
$\omega_\varepsilon$,
\begin{equation}
  K_X = -\langle\!\langle(1+q_XD)\,e_X(\bm\epsilon^\circ,\bm\epsilon)
  \rangle\!\rangle, \qquad
  K_{q_X} = -\langle\!\langle X D\,e_X(\bm\epsilon^\circ,\bm\epsilon)
  \rangle\!\rangle,
\end{equation}
for $X\in\{A,C,F,L,N\}$ and $e_X$ its coefficient above. Nothing in the
derivation depends on the form of the dispersion: for any rheology whose
moduli $\bc(\omega;\mathbf{p})$ are holomorphic in the lower half of the
complex plane, as causality requires, so that they may be evaluated at
$\omega_\varepsilon$, the kernel for a rheological parameter $p$ is
$K_p = -\langle\!\langle\partial_p\bc_\omega:(\nabla\sd_\omega\otimes
\nabla\bs_\omega)\rangle\!\rangle$: for a power-law $Q$, its exponent, or for
standard linear solids, their relaxation strengths and times. For isotropy,
$K_\kappa = -\langle\!\langle(1+q_\kappa D)\,\nabla\cdot\sd_\omega\,
\nabla\cdot\bs_\omega\rangle\!\rangle$ and $K_\mu = -\langle\!\langle
(1+q_\mu D)\,2\mathbf{d}^\circ:\mathbf{d}\rangle\!\rangle$ with $\mathbf{d}$
the deviatoric strain, as in \citet[][eqs~63--64]{LiuTromp2008}.
Attenuation kernels come at no extra cost. At fixed density the velocity
kernels follow from these: for isotropy, in relative perturbations,
$K_\alpha = 2(\kappa+\tfrac43\mu)K_\kappa$ and $K_\beta = 2\mu(K_\mu -
\tfrac43K_\kappa)$ with $K_\kappa$, $K_\mu$ the absolute kernels above
\citep[cf.][eq.~68]{LiuTromp2008}, and similarly for the five velocities of
transverse isotropy.

\paragraph{Density, at fixed moduli.}
For the spherically symmetric perturbations to which density is restricted,
the angular integral of the following is needed; we write it pointwise.
From \eqref{eq:S}, \eqref{eq:Ag}, \eqref{eq:B} and \eqref{eq:gradient},
\begin{equation}
\begin{aligned}
  K_\rho = \Big\langle\!\!\Big\langle\,&\omega_\varepsilon^2\,\sd_\omega
  \cdot\bs_\omega - \sd_\omega\cdot\nabla\phi_\omega - \bs_\omega\cdot
  \nabla\pd_\omega \\
  &- \bs_\omega\cdot\nabla\nabla\Phi\cdot\sd_\omega - (\bs_\omega\cdot
  \nabla\sd_\omega)\cdot\nabla\Phi + (\nabla\cdot\sd_\omega)(\bs_\omega\cdot
  \nabla\Phi)\Big\rangle\!\!\Big\rangle + \lambda. \label{eq:Krho}
\end{aligned}
\end{equation}
Every term is local; the equilibrium gravity's dependence on the density is
carried entirely by $\lambda$. For a spherically symmetric background the
gravitational terms simplify to
\begin{equation}
  \bs\cdot\nabla\nabla\Phi\cdot\sd + (\bs\cdot\nabla\sd)\cdot\nabla\Phi
  - (\nabla\cdot\sd)(\bs\cdot\nabla\Phi)
  = \left(4\pi G\rho - \frac{4g}{r}\right)s_rs^\circ_r
  + g\left(\bs_h\cdot\nabla_hs^\circ_r - s_r\nabla_h\cdot\sd_h\right),
\end{equation}
using $\nabla\nabla\Phi = \Phi''\br\br + (g/r)(\bI-\br\br)$ with
$\Phi'' = 4\pi G\rho - 2g/r$, and $(\bs\cdot\nabla)\br = \bs_h/r$, with
$\nabla_h$ the horizontal gradient. The coefficient $-4g/r$ is that of the
radial equations.

\paragraph{Density, at fixed velocities.}
Again for spherically symmetric perturbations: each modulus is proportional
to the density at fixed velocities, so $\rho K'_\rho = \rho K_\rho +
\sum_X XK_X$: for isotropy $\rho K_\rho +
\kappa K_\kappa + \mu K_\mu$ \citep[][eq.~68]{LiuTromp2008}, for transverse
isotropy $\rho K_\rho + AK_A + CK_C + FK_F + LK_L + NK_N$. The velocity
kernels follow as usual.

\paragraph{Source.}
For a moment tensor $\mathbf{M}_s$ at $\bx_s$, stepped in time,
$\langle\sd_\omega,\delta\mathbf{f}_\omega\rangle = \delta\mathbf{M}_s:
\bm\epsilon^\circ(\bx_s)/(i\omega_\varepsilon)$, giving the source kernels
by reciprocity at no extra cost.

\paragraph{At the source.}
The strain of a point moment tensor is not locally integrable: besides its
$R^{-3}$ near field, whose angular average vanishes, it has a part
concentrated at the source, as $\nabla\nabla R^{-1}$ has the local part of
Section~\ref{sec:lambda}. In the expansion in spherical harmonics this part
is explicit. At each degree the displacement jumps across $r=r_s$, by
$\mathbf{M}_s\cdot\br$ divided by $C$ or $L$ there, and its radial derivative
has a part $\jump{\bs}\,\delta(r-r_s)$, so that
\begin{equation}
  \bm\epsilon = \bm\epsilon_{\mathrm{reg}} + \sym\big(\br\otimes\jump{\bs}\big)
  \,\delta(r-r_s).
\end{equation}
The kernels of the moduli that multiply $\epsilon_{rr}$ and
$\bm\epsilon_{rh}$, namely $C$, $F$ and $L$, and their quality factors, have
accordingly a part concentrated on the sphere $r=r_s$; those of $A$, $N$ and
$\rho$ do not. For spherically symmetric perturbations it is a point term
$P_X\,\delta X(r_s)$ in $\delta J$, with
\begin{equation}
  P_X = -r_s^2\left\langle\!\!\left\langle(1+q_XD)\oint e_X\big(
  \bm\epsilon^\circ, \sym(\br\otimes\jump{\bs})\big)\dd\Omega
  \right\rangle\!\!\right\rangle,
\end{equation}
evaluated with the adjoint field at $r_s$, which is regular there unless the
receiver is at the depth of the source. Finite differences that perturb the
moduli around the source need it.

\paragraph{Discretisation.}
YSpec computes the spectra at $f_k = k\,\Delta f$ and forms the samples
$u_n = u(n\Delta t)$ by a fast Fourier transform. Discretised with the same
rules, $\langle\!\langle F\rangle\!\rangle$ becomes $2\Delta f\sum_k\Re
F(\omega_k)$ and the time integral in \eqref{eq:A} a sum over the samples,
so that $A_k = w_kR_k\sum_nj(t_n)\,e^{\varepsilon t_n}e^{2\pi if_kt_n}
\Delta t$. For these linear operations this coincides with the derivative of
the discrete misfit, with $\partial J/\partial u_n = j(t_n)\Delta t$, so the
kernels predict finite differences of the computed $J$ to the accuracy of
the solver.

\section{The adjoint potential}\label{sec:lambda}

\subsection{Strong form}

Where $\bM$ and $\vm$ are smooth, \eqref{eq:lambda} is
$-(4\pi G)^{-1}\nabla^2\lambda = \nabla\cdot\nabla\cdot\bM - \nabla\cdot\vm$,
and $\lambda$ is harmonic outside $V$ and decays at infinity. The source
contains double layers where $\bM_{nn}$ jumps (as $\rho$ does at
discontinuities), so $\lambda$ is itself discontinuous there, with
$\jump{\lambda} = -4\pi G\jump{\bM_{nn}}$; the weak form handles this
without comment. For the perturbations considered here only the spherical
average of $\lambda$ is needed, which reduces to the radial quadrature of
Section~\ref{sec:lambda1d}; the integral representation below is given for
comparison with \citet{Eaton2026}.

\subsection{Integral representation and comparison with
\citet{Eaton2026}}

Taking $\psi = -G/|\bx-\bx_0|$ in \eqref{eq:lambda}, for which
$\nabla^2\psi = 4\pi G\delta(\bx-\bx_0)$, gives
$\lambda(\bx_0) = -\int_V(\bM:\nabla\nabla\psi + \vm\cdot\nabla\psi)\dd V$.
Now $\nabla\psi = G\,\mathbf{R}/R^3$ with $\mathbf{R}=\bx-\bx_0$, and
\begin{equation}
  \nabla\nabla\psi = G\,\Pi + \frac{4\pi G}{3}\bI\,\delta(\bx-\bx_0),
  \qquad \Pi = \frac{\bI}{R^3} - \frac{3\mathbf{R}\otimes\mathbf{R}}{R^5},
\end{equation}
the integral over $\Pi$ being a principal value; the local term is needed
for $\nabla^2\psi = 4\pi G\delta$, since $\Pi$ is traceless. Hence
\begin{equation}
  \lambda(\bx_0) = -G\int_V\left[\bM:\Pi + \vm\cdot\frac{\bx-\bx_0}
  {|\bx-\bx_0|^3}\right]\dd V - \frac{4\pi G}{3}\,\mathrm{tr}\,\bM(\bx_0).
\end{equation}
With $\bM\leftrightarrow\tilde{\mathbf{K}}_\rho$ and $\vm\leftrightarrow
\tilde{\mathbf{k}}_\rho$ \citep[][eqs~23--24]{Eaton2026}, the first term is
the non-local part of their density kernel (their eq.~29 is for relative
perturbations, and so carries an extra factor $\rho$), so the two results
agree, given one remark: their eq.~A4 writes $\nabla\nabla\delta\Phi =
G\int\delta\rho\,\Pi$ without the local term, which is needed for
$\nabla^2\delta\Phi = 4\pi G\delta\rho$. Their numerical route, which takes
derivatives of solutions of Poisson's equation, includes it automatically.
Here a single scalar Poisson problem replaces their nine.

\subsection{One-dimensional perturbations}\label{sec:lambda1d}

For $\delta\rho(r)$ the perturbation of the equilibrium potential is radial
and explicit, $\delta\Phi' = 4\pi G\,\delta m(r)/r^2$ with $\delta m(r)
=\int_0^r\delta\rho\,r'^2\dd r'$, and $\delta\Phi'' = 4\pi G\delta\rho -
2\delta\Phi'/r$, so $\lambda$ need not be formed. With
$\nabla\nabla\delta\Phi = \delta\Phi''\br\br + (\delta\Phi'/r)
(\bI-\br\br)$ and the angular integrals
\begin{equation}
  \bar M_{rr}(r) = \oint\bM_{rr}\dd\Omega, \quad
  \bar M_h(r) = \oint(\bM_{\theta\theta}+\bM_{\phi\phi})\dd\Omega, \quad
  \bar m_r(r) = \oint\vm_r\dd\Omega,
\end{equation}
the gravitational part of the one-dimensional density kernel, defined by
$\delta J = \int_0^aK^{1D}_\rho\,\delta\rho\,r^2\dd r$, is
\begin{equation}
  K^{1D,\Phi}_\rho(r) = -4\pi G\left[\bar M_{rr}(r) + \int_r^a
  \left(\frac{\bar M_h - 2\bar M_{rr}}{r'} + \bar m_r\right)\dd r'\right],
\end{equation}
non-local only through an integral over the radii above. The local part is
the angular integral of the bracket in \eqref{eq:Krho}. YSpec computes $g$
from $\rho$, so finite differences in the density test this kernel
completely, gravitational part included.

\section{Why density is not perturbed laterally}\label{sec:lateral}

A background in hydrostatic equilibrium cannot remain so under a laterally
varying $\delta\rho$, or lateral topography: the initial stress acquires a
deviatoric part that the density does not determine, and the hydrostatic
equations no longer define the problem. One symptom: for a hydrostatic
background
\begin{equation}
  \nabla(\rho\bs\cdot\nabla\Phi) - \nabla\cdot(\rho\bs)\nabla\Phi =
  \rho\left[\nabla(\bs\cdot\nabla\Phi) - (\nabla\cdot\bs)\nabla\Phi\right]
\end{equation}
\citep[][eq.~B23]{AlAttarTromp2014}, the two sides differing by
$\nabla\rho\,(\bs\cdot\nabla\Phi) - (\nabla\rho\cdot\bs)\nabla\Phi$, which
vanishes when $\nabla\rho\parallel\nabla\Phi$. Perturbed by a laterally
varying $\delta\rho$ and the $\delta\Phi$ it produces, with $\nabla\Phi=g\br$
and $\nabla\rho = \rho'\br$, perturbing the left-hand form rather than the
right-hand one changes $\delta J$ by
\begin{equation}
  \left\langle\!\!\left\langle\int_V\left[\rho'\,\nabla_h\delta\Phi -
  g\,\nabla_h\delta\rho\right]\cdot\left(s_r\sd_h - s^\circ_r\bs_h\right)
  \dd V\right\rangle\!\!\right\rangle,
\end{equation}
where $\rho'$ includes $\jump{\rho}\delta_\Sigma$ at every discontinuity.
The three-dimensional density kernel therefore depends at first order on
which of two equivalent forms of the equations is perturbed; equivalently,
$A_g$ is symmetric only for a hydrostatic background
(Lemma~\ref{lem:symmetry}), and $A_g(\sd,\bs)$ and $A_g(\bs,\sd)$ have
different perturbations. \citet{LiuTromp2008}, and so \citet{Eaton2026},
perturb the left-hand form, and also omit the surface term
$-\int_\Sigma\jump{\rho(s^\circ_n\bs_h - s_n\sd_h)}\cdot\nabla_h\delta\Phi
\dd\Sigma$, which perturbing it generates on every boundary across which
$\rho$ jumps (it is the $\delta_\Sigma$ part of the expression above). None
of this arises for the perturbations considered here: at fixed density, and
for spherically symmetric $\delta\rho$, $\delta\Phi$ is radial and every term
vanishes.

\section{Topography on discontinuities}\label{sec:boundaries}

This section is an outline, to be completed, and is restricted to changes in
the radius of a discontinuity, the only topography that keeps the model
hydrostatic. Moving a discontinuity changes only the domains of the
integrals in \eqref{eq:L}. Mapping the perturbed
domain onto the reference one and using that every field satisfies its
equations and transmission conditions, the first variation reduces to the
jump across the boundary of the normal--normal component of the
energy--momentum tensor of the Lagrangian density,
$\mathbb{E} = \ell\bI - \sum_v(\nabla v)^{\mathsf T}\cdot\partial\ell/
\partial(\nabla v)$, summed over the fields $v$ on which $\ell$ depends
through first derivatives, with an additional tangential term on
$\Sigma_{FS}$, where the boundary may move relative to the slipping fluid.
For the elastic and inertial terms, $\ell = \rho\omega^2\sd\cdot\bs -
\nabla\sd:\bc:\nabla\bs$, and $\bn\cdot\mathbb{E}\cdot\bn = \ell +
(\bn\cdot\bT)\cdot\partial_n\sd + (\bn\cdot\bT^\circ)\cdot\partial_n\bs$,
which corresponds, up to their sign convention and the time domain, to the
kernel of \citet[][eq.~60]{LiuTromp2008}. The constraint
contributes $\jump{\rho\lambda - (4\pi G)^{-1}g\,\partial_n\lambda}$ at a
spherical boundary, and the gravitational terms of $A_g$ involve
$\nabla\nabla\Phi$, so that their energy--momentum tensor needs the
second-order form. The kernel is one number per discontinuity, and finite
differences that move a discontinuity will check it.

\section{Evaluation in YSpec}\label{sec:yspec}

For a spherically symmetric background, YSpec gives the forward field of
each source and the field of a point force at each receiver at every radius
wanted, in the frame with the source (or receiver) at the pole, as
coefficients of generalised spherical harmonics with $|m|\le2$ (or $1$), and
their strains as the spin components $N=-2,\dots,2$.
\begin{itemize}
\item One-dimensional kernels need only angular integrals of products of the
  two fields. The adjoint field is expanded in the frame of the source, where
  a force away from the pole excites every order but only $|m|\le2$ meet the
  forward field: its weights are those of the force at the pole with
  $X^N_{lm}(0)$ replaced by $X^N_{lm}(\Delta)\,e^{-im\phi_r}$. By
  orthogonality, $\oint Y^N_{lm}Y^{-N}_{l,-m}\dd\Omega = (-1)^{N+m}$, each
  integral is then a sum over degree and order of products of radial
  functions, and for displacements $\oint\bs\cdot\sd\dd\Omega =
  \sum_{lm}(-1)^m[U_mU^\circ_{-m} + \zeta^2(V_mV^\circ_{-m} +
  W_mW^\circ_{-m})]$; there is no spatial mesh. The kernels are evaluated at
  Gauss--Legendre nodes in elements between the knots of the model, the
  source and the receiver, and the radial integral in $K^{1D,\Phi}_\rho$ is
  taken on the same nodes. With $\bar M_h - 2\bar M_{rr} + r\bar m_r =
  \rho\big[\zeta^2\sum(-1)^m(V_mU^\circ_{-m} + U_mV^\circ_{-m}) -
  4\sum(-1)^mU_mU^\circ_{-m}\big]$ summed over degree, its integrand is the
  same combination as the gravitational term $g(\bs_h\cdot\nabla_hs^\circ_r
  - s_r\nabla_h\cdot\sd_h) - (4g/r)s_rs^\circ_r$ of the local kernel, divided
  by $g/\rho$.
\item Three-dimensional kernels, for the elastic parameters, need the
  fields on a mesh. Expansions of spin-weighted fields rotate one $N$ at a
  time, so each field is rotated from its own frame to geographic
  coordinates with Wigner matrices (cheaply, since only a few orders are
  present) and synthesised on a Gauss--Legendre grid at each radius; the
  kernels are pointwise products integrated over frequency. No $\lambda$ is
  needed. A cross-section through source and receiver is far cheaper than a
  full mesh.
\item Checks: finite differences of $J$ for spherically symmetric
  perturbations of every parameter, density with gravity included, and the
  angular integrals of the three-dimensional elastic kernels against the
  one-dimensional ones.
\end{itemize}
One restriction: a receiver must be where YSpec can place a source, that is,
in a solid.

\begin{thebibliography}{9}
\bibitem[Adourian et~al.(2024)]{Adourian2024}
Adourian, S., Dursun, M.S., Lau, H.C.P. \& Al-Attar, D., 2024. Adjoint
sensitivity kernels for free oscillation spectra, \emph{Geophys. J. Int.},
238, 257--271.
\bibitem[Al-Attar \& Tromp(2014)]{AlAttarTromp2014}
Al-Attar, D. \& Tromp, J., 2014. Sensitivity kernels for viscoelastic
loading based on adjoint methods, \emph{Geophys. J. Int.}, 196, 34--77.
\bibitem[Dahlen et~al.(2000)]{Dahlen2000}
Dahlen, F.A., Hung, S.-H. \& Nolet, G., 2000. Fr\'echet kernels for
finite-frequency traveltimes---I. Theory, \emph{Geophys. J. Int.}, 141,
157--174.
\bibitem[Dahlen \& Tromp(1998)]{DahlenTromp1998}
Dahlen, F.A. \& Tromp, J., 1998. \emph{Theoretical Global Seismology},
Princeton University Press.
\bibitem[Eaton et~al.(2026)]{Eaton2026}
Eaton, W., Gharti, H.N. \& Tromp, J., 2026. Sensitivity of prompt
elasto-gravitational signals to 3-D Earth structure, \emph{Geophys. J.
Int.}, 245, 1--25.
\bibitem[Liu \& Tromp(2008)]{LiuTromp2008}
Liu, Q. \& Tromp, J., 2008. Finite-frequency sensitivity kernels for global
seismic wave propagation based upon adjoint methods, \emph{Geophys. J.
Int.}, 174, 265--286.
\end{thebibliography}

\end{document}
